\documentclass[a4paper, serif, 10pt]{moderncv}

%% moderncv
% style and color
\moderncvstyle{banking}
\moderncvcolor{black}

%% page layout/margins
\usepackage[top=2.5cm, bottom=2.5cm, left=1.5cm, right=1.5cm]{geometry}

\nopagenumbers{}

%% Babel config for Brazilian Portuguese language
% ref:
% - https://latex3.github.io/babel/guides/locale-portuguese.html
\usepackage[brazilian]{babel}

%% fontspec
\usepackage{fontspec}

\IfFontExistsTF{Linux Libertine}{
  \setmainfont[Ligatures={TeX, Common}, Numbers=OldStyle]{Linux Libertine}
}{}
\IfFontExistsTF{Linux Libertine O}{
  \setmainfont[Ligatures={TeX, Common}, Numbers=OldStyle]{Linux Libertine O}
}{}

%% CTeX
% CJK support, used to show CJK characters in the resume
%
% - fontset=none: disable builtin fontset but instead set the CJK font manually
% - heading=false: disable ctex heading
% - punct=kaiming: use kaiming punctuations styles for CJK
% - scheme=plain: use plain scheme, do not override \normalsize font size
% - space=auto: space settings for CJK characters
%
% ref:
% - http://ctan.mirrorcatalogs.com/language/chinese/ctex/ctex.pdf
\usepackage[UTF8, heading=false, punct=kaiming, scheme=plain, space=auto]{ctex}

\IfFontExistsTF{Noto Serif CJK SC}{
  \setCJKmainfont{Noto Serif CJK SC}
}{}
\IfFontExistsTF{Noto Sans CJK SC}{
  \setCJKsansfont{Noto Sans CJK SC}
}{}

%% line spacing
\usepackage{setspace}
\setstretch{1.125}

%% URL styling - use normal text instead of monospace
\urlstyle{same}

% Auto-underline all links
% Defer until \begin{document} so hyperref is already loaded
\AtBeginDocument{%
  \let\oldhref\href
  \renewcommand{\href}[2]{\oldhref{#1}{\underline{#2}}}%
}

%% Patch \httplink and \httpslink to handle full URLs
% The original moderncv macros always prepend http:// or https:// to URLs.
% This causes issues when the URL already has a protocol (e.g., https://example.com)
% resulting in malformed URLs like https://https://example.com.
% This patch checks if the URL already contains "://" and uses it directly if so.
% Note: We use \str_set:Nx (x-expansion) to expand macros like \@homepage first.
\ExplSyntaxOn
\str_new:N \g_yamlresume_url_str

\RenewDocumentCommand{\httpslink}{O{}m}{
  \str_set:Nx \g_yamlresume_url_str {#2}
  \str_if_in:NnTF \g_yamlresume_url_str {://}
    { \url{#2} }
    { \url{https://#2} }
}

\RenewDocumentCommand{\httplink}{O{}m}{
  \str_set:Nx \g_yamlresume_url_str {#2}
  \str_if_in:NnTF \g_yamlresume_url_str {://}
    { \url{#2} }
    { \url{http://#2} }
}
\ExplSyntaxOff

\name{Rafael Nogueira Almeida}{}
\title{Engenheiro DevOps Sênior | Cloud, Kubernetes e Platform Engineering}
\phone[mobile]{+55 11 98765-4321}
\email{rafael.nogueira.almeida@outlook.com}
\homepage{https://www.rafaelalmeida.dev}

\address{Rua Tabapuã, 1234, Apto 82 - Itaim Bibi, São Paulo, São Paulo, Brasil, 04533-003}{}{}

\extrainfo{{\small \faGithub}\ \href{https://github.com/rafaelnalmeida}{@rafaelnalmeida} {} {} {} • {} {} {} 
{\small \faLinkedin}\ \href{https://www.linkedin.com/in/rafael-almeida-devops}{@rafael-almeida-devops} {} {} {} • {} {} {} 
{\small \faGitlab}\ \href{https://gitlab.com/rafaelnalmeida}{@rafaelnalmeida}}

\begin{document}

\maketitle

\section{Informações Básicas}

\cvline{}{Engenheiro DevOps com mais de 8 anos de experiência em nuvem, automação de infraestrutura e plataformas de entrega contínua.
%
\begin{itemize}
\item Especialista em Kubernetes, Terraform e AWS, com foco em confiabilidade, escalabilidade e eficiência de custos
\item Reduzi em 70\% o tempo de provisionamento de ambientes e em 60\% o lead time de deploy em times de produto
\item Vivência em práticas SRE, definição de SLOs, resposta a incidentes e cultura de pós-mortem sem culpa
\item Atuação em mentoria, documentação técnica e disseminação de boas práticas de Platform Engineering
\end{itemize}}
\section{Formação}

\cventry{mar. de 2013–dez. de 2017}
        {Bacharelado, Engenharia de Computação, Nota: 8.4/10}
        {Universidade Estadual de Campinas}
        {\url{https://www.unicamp.br}}
        {}
        {\begin{itemize}
\item Base sólida em sistemas distribuídos, redes e administração de sistemas Linux
\item Projeto de conclusão sobre orquestração de contêineres e escalabilidade horizontal de serviços
\item Monitoria voluntária na disciplina de Sistemas Operacionais
\end{itemize}
\textbf{Cursos}: Sistemas Operacionais, Redes de Computadores, Sistemas Distribuídos, Banco de Dados, Engenharia de Software, Segurança da Informação, Computação em Nuvem}

\cventry{fev. de 2021–nov. de 2022}
        {Diploma, Arquitetura de Soluções em Cloud Computing}
        {FIAP - Faculdade de Informática e Administração Paulista}
        {\url{https://www.fiap.com.br}}
        {}
        {\begin{itemize}
\item Pós-graduação lato sensu com foco em arquitetura de nuvem, multi-cloud e governança
\item Trabalho final sobre plataformas internas de autoatendimento para times de engenharia
\end{itemize}}
\section{Experiência Profissional}

\cventry{mar. de 2022–Atual}
        {Engenheiro DevOps Sênior}
        {Nuvem Viva Tecnologia}
        {\url{https://www.nuvemviva.com.br}}
        {}
        {\begin{itemize}
\item Lidero a modernização da plataforma de microsserviços em Kubernetes (Amazon EKS), atendendo mais de 40 squads de produto
\item Estruturei a infraestrutura como código com Terraform e Terragrunt, reduzindo em 70\% o tempo de provisionamento de novos ambientes
\item Implementei pipelines de entrega contínua com GitLab CI e Argo CD, diminuindo o lead time de deploy de 2 dias para 40 minutos
\item Defini SLOs, alertas e runbooks em Prometheus e Grafana, reduzindo o MTTR de incidentes críticos em 55\%
\item Conduzi iniciativas de FinOps que geraram economia anual de R\$ 1,2 milhão em custos de nuvem
\item Mentorei 5 engenheiros e engenheiras em práticas SRE, segurança de contêineres e análise de incidentes
\end{itemize}
\textbf{Palavras-chave}: Kubernetes, Terraform, AWS, SRE, FinOps}

\cventry{jan. de 2020–fev. de 2022}
        {Engenheiro DevOps}
        {PagFácil Pagamentos}
        {\url{https://www.pagfacil.com.br}}
        {}
        {\begin{itemize}
\item Migrei cargas monolíticas para contêineres Docker orquestrados em Kubernetes, elevando a disponibilidade de 99,5\% para 99,95\%
\item Automatizei o provisionamento de servidores com Ansible e Packer, eliminando tarefas manuais repetitivas de operação
\item Implementei a stack de observabilidade com ELK e Datadog, cobrindo 100\% dos serviços críticos de pagamento
\item Introduzi gestão de segredos com Vault e políticas de menor privilégio alinhadas ao PCI DSS
\item Atuei em plantões on-call e na resposta a incidentes em ambiente de alto volume transacional
\end{itemize}
\textbf{Palavras-chave}: CI/CD, Docker, Jenkins, Observabilidade, Vault}

\cventry{jul. de 2017–dez. de 2019}
        {Analista de Infraestrutura e Automação}
        {Sistemas Origem Consultoria}
        {}
        {}
        {\begin{itemize}
\item Administrei ambientes Linux (Debian e RHEL) e clusters de virtualização VMware para clientes de médio porte
\item Desenvolvi scripts em Python e Bash para automação de backups, monitoramento e deploys
\item Configurei balanceadores Nginx, certificados TLS e VPNs site-to-site entre unidades dos clientes
\item Documentei procedimentos operacionais e treinei equipes de suporte de primeiro e segundo nível
\end{itemize}
\textbf{Palavras-chave}: Linux, Ansible, Virtualização, Scripts, Nginx}
\section{Idiomas}

\cvline{Português}{Nativo ou Bilíngue}
\cvline{Inglês}{Proficiência Profissional Fluente \hfill \textbf{Palavras-chave}: TOEFL iBT 108, Reuniões técnicas diárias}
\cvline{Espanhol}{Proficiência Profissional Limitada \hfill \textbf{Palavras-chave}: Leitura de documentação}
\section{Competências}

\cvline{Cloud \& Infraestrutura como Código}{Especialista \hfill \textbf{Palavras-chave}: AWS, Azure, Terraform, Terragrunt, Ansible, Packer}
\cvline{Contêineres \& Orquestração}{Especialista \hfill \textbf{Palavras-chave}: Docker, Kubernetes, Helm, Argo CD, Istio}
\cvline{CI/CD \& Automação}{Avançado \hfill \textbf{Palavras-chave}: GitLab CI, GitHub Actions, Jenkins, Python, Bash}
\cvline{Observabilidade \& SRE}{Avançado \hfill \textbf{Palavras-chave}: Prometheus, Grafana, Datadog, OpenTelemetry, ELK}
\cvline{Segurança \& Conformidade}{Intermediário \hfill \textbf{Palavras-chave}: Vault, IAM, Trivy, PCI DSS, LGPD}
\section{Prêmios}

\cventry{dez. de 2023}
        {Nuvem Viva Tecnologia}
        {Prêmio Inovação Interna - Plataforma de Autoatendimento}
        {}
        {}
        {Reconhecimento pelo desenvolvimento da plataforma interna de provisionamento de ambientes, que reduziu em 70\% o tempo de entrega de infraestrutura para as squads de produto.}

\cventry{jan. de 2021}
        {PagFácil Pagamentos}
        {Destaque Técnico do Ano}
        {}
        {}
        {Premiação anual concedida aos profissionais de maior impacto técnico, pela liderança da migração para contêineres e pela melhoria dos indicadores de disponibilidade da plataforma de pagamentos.}
\section{Certificados}

\cventry{mar. de 2023}
        {Amazon Web Services}
        {AWS Certified Solutions Architect - Professional}
        {\url{https://aws.amazon.com/certification/}}
        {}
        {}

\cventry{jul. de 2022}
        {Cloud Native Computing Foundation}
        {Certified Kubernetes Administrator (CKA)}
        {\url{https://www.cncf.io/certifications/cka/}}
        {}
        {}

\cventry{nov. de 2021}
        {HashiCorp}
        {HashiCorp Certified: Terraform Associate}
        {\url{https://www.hashicorp.com/certifications/terraform-associate}}
        {}
        {}
\section{Publicações}

\cventry{ago. de 2023}
        {Escalando clusters Kubernetes com Karpenter no Amazon EKS}
        {Blog Técnico Nuvem Viva}
        {\url{https://medium.com/@rafaelnalmeida/escalando-clusters-kubernetes-com-karpenter-no-amazon-eks-9f1c2b7a4d13}}
        {}
        {\begin{itemize}
\item Artigo técnico sobre provisionamento dinâmico de nós e redução de custos em clusters EKS
\item Comparativo entre Cluster Autoscaler e Karpenter em cargas de trabalho variáveis
\item Apresenta resultados reais com economia de 30\% em custos de computação
\end{itemize}}
\section{Referências}

\cventry{\emaillink[camila.santos@nuvemviva.com.br]{camila.santos@nuvemviva.com.br}}
        {Gerente de Plataforma na Nuvem Viva Tecnologia}
        {Camila Ferreira dos Santos}
        {+55 11 97654-3210}
        {}
        {Rafael é referência técnica em Kubernetes e Terraform no time. Conduz iniciativas complexas com autonomia, comunica-se bem com áreas de produto e eleva o nível técnico de quem trabalha ao seu lado.}

\cventry{\emaillink[eduardo.nakamura@pagfacil.com.br]{eduardo.nakamura@pagfacil.com.br}}
        {Diretor de Engenharia na PagFácil Pagamentos}
        {Eduardo Nakamura}
        {+55 11 96543-2109}
        {}
        {Durante dois anos, Rafael foi responsável por nossa plataforma de entrega contínua. Sua atuação reduziu drasticamente o tempo de deploy e trouxe uma cultura de confiabilidade que permanece até hoje.}
\section{Projetos}

\cventry{mar. de 2023–nov. de 2023}
        {Plataforma interna que permite às squads provisionar ambientes completos por meio de pull requests.}
        {Plataforma de Autoatendimento de Infraestrutura}
        {\url{https://github.com/rafaelnalmeida/plataforma-autoatendimento}}
        {}
        {\begin{itemize}
\item Construída com Backstage, Terraform, Argo CD e Kubernetes em modelo GitOps
\item Reduziu de 3 dias para 25 minutos o tempo de criação de ambientes de homologação
\item Adotada por 32 squads, com mais de 400 ambientes provisionados no primeiro semestre
\end{itemize}
\textbf{Palavras-chave}: Backstage, Terraform, Argo CD, Kubernetes, GitOps}

\cventry{fev. de 2021–out. de 2021}
        {Stack de métricas, logs e traces para monitorar transações críticas 24 horas por dia.}
        {Observabilidade para Pagamentos em Tempo Real}
        {}
        {}
        {\begin{itemize}
\item Instrumentação de serviços com OpenTelemetry e correlação de traces ponta a ponta
\item Dashboards e alertas baseados em SLOs de latência e taxa de erro
\item Redução de 55\% no tempo médio de diagnóstico de falhas em produção
\end{itemize}
\textbf{Palavras-chave}: OpenTelemetry, Grafana, Datadog, SLO}
\section{Interesses}

\cvline{Esportes}{Corrida de rua, Ciclismo, Natação}
\cvline{Comunidade Técnica}{Palestras, Open Source, Mentoria}
\cvline{Música}{Guitarra, Vinis}
\section{Voluntariado}

\cventry{mar. de 2019–nov. de 2023}
        {Voluntário na Organização}
        {DevOpsDays São Paulo}
        {\url{https://devopsdays.org/sao-paulo}}
        {}
        {\begin{itemize}
\item Participei da organização de edições do DevOpsDays São Paulo, apoiando curadoria de palestras e logística do evento
\item Ministrei workshops introdutórios de Kubernetes e Terraform para mais de 200 participantes
\item Atuei como mentor em programas de entrada de novos profissionais nas áreas de DevOps e SRE
\end{itemize}}
    
\end{document}